\section{Introducción}
Una tasa positiva de crecimiento posterior a una recesión no implica recuperar el nivel de actividad previo. Esta distinción es relevante al comparar economías regionales con estructuras productivas diferentes. Un departamento dependiente de un proyecto extractivo puede registrar variaciones pronunciadas de producción por cambios en capacidad o interrupciones operativas; otro con mayor participación de servicios puede mostrar tasas más estables. La comparación de tasas anuales ofrece por ello una imagen distinta de la comparación de niveles normalizados.

Este documento examina cómo se relaciona la estructura productiva previa a 2020 con la volatilidad del crecimiento y la recuperación observada hasta \LevelEnd. La pregunta empírica es si los perfiles productivos presentan cambios diferentes durante el choque de \ShockYear\ y el periodo \PostStart--\PostEnd. El análisis separa tres operaciones: construir una clasificación con información previa, describir la dispersión histórica y estimar cambios relativos dentro de cada departamento.

La contribución es integrar esas operaciones en un flujo reproducible. Tratamos las participaciones sectoriales como composiciones que suman 100\%, en lugar de interpretarlas como variables libres e independientes. Comparamos cuatro familias de agrupamiento y tres representaciones de los datos. Finalmente, contrastamos una misma estructura de interacciones en modelos cuya dependiente es el crecimiento anual o el nivel de VAB real indexado.

El principal resultado es que el perfil minero/enclave permanece por debajo de su nivel medio de 2019 durante 2021--2023, mientras los demás perfiles terminan por encima de ese umbral. El modelo en índices estima un cambio relativo negativo para ese grupo. Su interpretación requiere cautela: hay solamente \EnclaveCount\ departamentos enclave y la clasificación utiliza una reducción de dimensión que no conserva toda la geometría composicional. La evidencia describe una recuperación desigual; no permite atribuirla exclusivamente a la pandemia, a conflictos o a una ausencia de diversificación.

\section{Datos, unidad de medida y muestras}
\subsection{Fuente y armonización}
Utilizamos los 24 libros departamentales del INEI conservados en el repositorio. El cuadro 1 contiene el \textit{valor agregado bruto} a precios constantes de 2007 y el cuadro 2 su estructura porcentual por actividades. Las notas de los libros consultados señalan información disponible a junio de 2026. El portal del INEI ofrece el marco de las cuentas departamentales y publicaciones de referencia \cite{inei}; la versión exacta utilizada se identifica mediante los archivos locales y sus huellas SHA-256, sin suponer que una publicación enlazada reproduce la misma revisión.

La columna heredada \texttt{PIB\_total} de la base consolidada corresponde a la fila \textit{Valor Agregado Bruto}. La salida final utiliza \texttt{VAB\_total} y \texttt{VAB\_idx2019}. Esta distinción evita confundir VAB con PBI, que también incorpora impuestos netos sobre los productos. Las tasas y los índices calculados se refieren a la misma magnitud de VAB real.

La base original presenta 12 columnas sectoriales porque una actividad residual aparece bajo dos nombres según el departamento. Unificamos \textit{Otros Servicios} y \textit{Otros Actividades/1}; resultan \Sectors\ actividades comparables. La categoría residual incluye otros servicios y, según las notas de los libros, pesca y acuicultura. Por ello no debe interpretarse como una medición exclusiva del sector servicios. Verificamos el cierre de las participaciones a 100\% antes de calcular cocientes.

La unidad transversal es la consignada en cada libro departamental; se conserva la cobertura de Lima tal como figura en su archivo. No se construye una observación adicional para Callao ni se mezclan subunidades de Lima con el total del libro. Las medias por perfil otorgan el mismo peso a cada departamento: no representan el VAB agregado nacional ni un promedio ponderado por tamaño económico.

\subsection{Ventanas con propósitos diferentes}
La base local contiene \BaseObs\ observaciones entre \BaseStart\ y \BaseEnd. Los años 2008--2023 están completos para los \Departments\ departamentos; 2024 y 2025 cuentan con 22 unidades cada uno. Se conservan esos años en la base, pero se excluyen de las muestras de estimación para evitar que la cobertura incompleta altere las comparaciones.

La clasificación utiliza el promedio aritmético de participaciones de \PreStart--\PreEnd. Esta ventana precede al COVID-19 y evita clasificar departamentos con los propios resultados de la recuperación. No garantiza exogeneidad estricta: la composición puede relacionarse con tendencias previas, inversión o instituciones no observadas.

Los modelos de crecimiento utilizan \RateStart--\RateEnd, con \RatesObs\ observaciones balanceadas. La variante M4 retira una observación para sensibilidad y utiliza \NFour\ filas. Los modelos de niveles utilizan \LevelStart--\LevelEnd, con \LevelsObs\ observaciones y siete años por departamento. El año base es \BaseYear. Las tasas disponibles en la base consolidada están redondeadas a dos decimales; se verifican contra los niveles anuales disponibles.
